\begin{pastexam}[label={px:ust-113-ee-ds-07}]{臺聯大電機類 113 學年度 第 7 題}
（單選；每題 3 分，答錯不倒扣）

A sorting algorithm is considered stable if:
\begin{enumerate}[label=(\Alph*)]
\item It requires O($n^2$) extra space for sorting.
\item It employs a divide and conquer strategy for sorting.
\item It maintains the relative order of equal elements as they appeared in the original array.
\item It sorts the elements in O(log n) time complexity.
\end{enumerate}
\end{pastexam}

\begin{pastsoln}
答案是 (C)。

穩定的定義：鍵相同的兩筆資料，排序後的先後順序和排序前一樣（資結書第~14.1~節）。
\begin{itemize}
\item (A) 講的是額外空間，和穩定無關；而且常見的排序沒有一個需要 $O(n^2)$ 額外空間。 錯。
\item (B) 分治是設計方法。 合併排序是分治而且穩定，快速排序也是分治卻不穩定，兩者無關。 錯。
\item (C) 正是定義。 對。
\item (D) 講的是時間；任何排序至少要看過每個元素，不可能 $O(\log n)$。 錯。
\end{itemize}
\end{pastsoln}
